%! Trigonometry and Polar Coordinates — Unit 3, Lesson 3.13 — Exit Ticket
\documentclass[10pt]{article}
\usepackage{precalculus-article}
\usepackage{precalculus-boxes}
\newcommand{\TallMath}[1]{$\displaystyle #1\rule[-1.4em]{0pt}{3.2em}$}

\begin{document}

\pageheader{Unit 3, Lesson 3.13}{Exit Ticket}
\namedateperiod

\vspace{0.12in}

\begin{notesbox}{Exit Ticket --- Answer independently}

\textbf{1.}\quad Convert the polar coordinates $(4,\, \pi/3)$ to rectangular coordinates.

\vspace{6pt}
$x = 4\cos\!\left(\dfrac{\pi}{3}\right) = $ \underline{\hspace{3.5cm}}

\vspace{8pt}
$y = 4\sin\!\left(\dfrac{\pi}{3}\right) = $ \underline{\hspace{3.5cm}}

\vspace{6pt}
Rectangular form: $(\underline{\hspace{2cm}},\; \underline{\hspace{2cm}})$

\vspace{0.18in}
\textbf{2.}\quad Convert the rectangular coordinates $(-3,\, 3)$ to polar form
with $r > 0$ and $0 \le \theta < 2\pi$.

\vspace{6pt}
$r = \sqrt{(-3)^2 + 3^2} = $ \underline{\hspace{3cm}}

\vspace{8pt}
$\theta = $ \underline{\hspace{5cm}} (show quadrant reasoning)

\vspace{6pt}
Polar form: $(\underline{\hspace{2cm}},\; \underline{\hspace{2cm}})$

\vspace{0.18in}
\textbf{3.}\quad Write the complex number $2 - 2i$ in polar form $(r\cos\theta) + i(r\sin\theta)$.

\vspace{6pt}
$r = $ \underline{\hspace{2.5cm}} \qquad $\theta = $ \underline{\hspace{2.5cm}}

\vspace{6pt}
Polar form: \underline{\hspace{7cm}}

\end{notesbox}

\end{document}
